\documentclass[12pt]{ca-capstone}
\geometry{margin=25mm}
\usepackage{setspace}
\setstretch{1.1}
\capstonerefs{knuth-taocp2,crandall-pomerance,welschenbach}
\begin{document}
\capstonetitle{Быстрое умножение больших чисел}
\begin{center}\textbf{Теоретическая справка для преподавателя. Не выдавать студентам.}\end{center}

\section{Свёртка и представление чисел}
Запишем $a=\sum_{i} a_{i}B^{i}$, $b=\sum_{j} b_{j}B^{j}$, $0\le a_{i},b_{j}<B$. Коэффициенты
произведения до нормализации переносов образуют линейную свёртку
\[
  c_{k}=\sum_{i+j=k}a_{i}b_{j},\qquad 0\le k<|a|+|b|-1.
\]
Циклическая свёртка длины $N$ отличается индексацией $i+j\pmod N$. Линейную свёртку получают
дополнением нулями до $N\ge |a|+|b|-1$. Коэффициенты $c_k$ не являются цифрами основания $B$:
их нельзя объединять побитовым <<или>>; нужны сложения $c_{k}B^{k}$ с учётом переносов.

\section{Карацуба и Тоом-3}
Для $a=a_1B^m+a_0$, $b=b_1B^m+b_0$ метод Карацубы использует
\[
 ab=z_2B^{2m}+(z_3-z_2-z_0)B^m+z_0,\quad
 z_0=a_0b_0,\ z_2=a_1b_1,\ z_3=(a_0+a_1)(b_0+b_1).
\]
Три умножения вместо четырёх дают сложность $O(n^{\log_2 3})$. Рекурсия должна иметь явный
базовый случай; неравная длина обрабатывается выравниванием по большему числу.

В Тоом-3 числа разбиваются на три блока и рассматриваются как значения квадратичных многочленов.
Типичные точки: $0,1,-1,-2,\infty$. Пять произведений значений и интерполяция восстанавливают
коэффициенты произведения; сложность $O(n^{\log_3 5})$. В эталоне используются формулы
\[
\begin{aligned}
r_0&=v_0,& r_4&=v_\infty,& r_3&=(v_{-2}-v_1)/3,\\
r_1&=(v_1-v_{-1})/2,& r_2&=v_{-1}-v_0,&&\\
r_3&=(r_2-r_3)/2+2r_4,&r_2&=r_2+r_1-r_4,&r_1&=r_1-r_3.
\end{aligned}
\]
Деления должны быть точными. Значения в отрицательных точках могут быть отрицательными, поэтому
рекурсия обязана поддерживать знак; базовый школьный метод работает только с абсолютными величинами.

\section{БПФ в конечном поле}
Пусть $N$ --степень двойки, $p=kN+1$ --простое, а $\omega$ имеет точный порядок $N$ в
$\GF{p}^*$. Прямое и обратное преобразования:
\[
 A_t=\sum_{j=0}^{N-1}a_j\omega^{tj}\pmod p,\qquad
 a_j=N^{-1}\sum_{t=0}^{N-1}A_t\omega^{-tj}\pmod p.
\]
Свёртка равна обратному преобразованию поточечного произведения преобразований. Итеративное
radix-2 БПФ требует перестановки битов индексов и $O(N\log N)$ модульных операций; обратное
получается заменой $\omega$ на $\omega^{-1}$ и умножением на $N^{-1}$.

Для автоматического выбора параметров достаточно:
\begin{enumerate}
  \item выбрать минимальную степень двойки $N$ не меньше требуемой длины;
  \item оценить $\max c_k\le N(B-1)^2=C$;
  \item перебирать $k$, пока $p=kN+1>C$ не окажется простым;
  \item для $h=2,3,\ldots$ вычислять $\omega=h^k\bmod p$, пока
    $\omega^{N/2}\ne1\pmod p$. Поскольку порядок делит степень двойки $N$, это гарантирует
    точный порядок $N$.
\end{enumerate}
Условие $p>C$ позволяет восстановить неотрицательные коэффициенты однозначно. Выбор только
простого $p\equiv1\pmod N$ недостаточен: без оценки $C$ результат известен лишь по модулю $p$.

\section{Сложность и типичные ошибки}
Школьный метод требует $O(n^2)$ поразрядных операций. Стоимость преобразований доминирует в
БПФ-умножении, но подбор простого и большие константы влияют на точку пересечения с рекурсивными
методами. Проверяются: точность интерполяционных делений, знак промежуточных значений, дополнение
до линейной длины, точный (а не просто делящий $N$) порядок $\omega$, отдельное масштабирование
на $N^{-1}$ и нормализация коэффициентов после обратной свёртки.

\clearpage
\section{Контрольный список}
Эталонные тесты связывают каждый метод с встроенным умножением; свёртку --с двойным циклом.
Отдельные группы: нули и знаки; $B-1$ во всех цифрах; сильно неравная длина; переходы $C$ через
$p$; неверный корень; $N=1$; проверка $A_t$ прямым ДПФ. Benchmark помечает метод и форму входа,
поэтому медленный класс можно отличить от ошибки алгоритма.

\references
\end{document}
